\section{伦理与公平数据：从研究成果到数据治理}

很多 ``Autonomous Agent'' 讨论只聚焦策略最优，但 Stone 在中段专门提到 ethical AI initiative 和 Phoebe 数据集，说明他将伦理治理视为 Agent 系统的一部分，而非部署后补丁。

\subsection{为什么伦理问题在 Agent 里更尖锐}
当系统从 ``回答问题'' 走向 ``持续决策''，伦理风险会被放大。Agent 有状态、可记忆、可长期行动，它造成的偏差不是单条输出，而是策略级行为偏置。

\begin{importantbox}{课程中的关键观点}
\textbf{课堂提示：}Stone 明确强调了两件事：第一，数据集不只看规模，还要看构建过程是否合规；第二，公平评测应面向全球多样人群，而不是局限于单地区、单语种分布。
\textit{``first globally diverse consensually collected fairness evaluation dataset''}（约 00:17:07）。
\end{importantbox}

\begin{knowledgebox}{Phoebe 的方法学意义}
课程提到 Phoebe（Nature 2025）时，重点并不在 ``10,000 图像'' 这个数量，而在以下流程要点：
\begin{itemize}
    \item 先识别数据构建过程中的伦理与法律问题，再定采集协议；
    \item 强调 \textbf{consensually collected}，避免 ``抓取即合规'' 的默认假设；
    \item 用数据集去评测模型偏见，而不是只做数据陈列。
\end{itemize}
\end{knowledgebox}

\subsection{把 fairness 变成系统接口，而不是论文附录}
这部分对 LLM Agent 尤其重要。很多团队在上线前只做一次性安全评审，忽略了上线后策略漂移。Stone 这段内容给出的启发是：\textbf{应把公平评测做成持续流程与平台能力}。

\begin{table}[H]
\centering
\begin{tabular}{p{3.2cm}p{4.6cm}p{4.2cm}}
\toprule
\textbf{治理环节} & \textbf{课程线索} & \textbf{工程落地点} \\
\midrule
数据采集 & consensually collected & 用户授权记录、数据来源台账 \\
标注与描述 & ethical/legal issues first & 敏感属性审查、标注规范审计 \\
评测目标 & evaluate bias in AI models & 分人群分任务指标看板 \\
发布与复现 & establish industry standard & 版本化评测基线与复现实验脚本 \\
\bottomrule
\end{tabular}
\caption{从课程观点到公平治理流水线}
\end{table}

\begin{warningbox}{常见短板：把公平只当成离线分数}
如果公平只在离线 benchmark 上看一次，部署后很快会失效。Agent 与真实用户交互后，输入分布和行为路径都会变。需要持续监测与回归测试，而不是 ``一次过检''。
\end{warningbox}

\subsection{本章小结}
Stone 讨论 Phoebe 的意义在于给出一个 ``可执行'' 的伦理路线：先治理流程，再谈模型分数；先建立标准，再扩展能力。

